% Methodological note of pyxtdpthresh: Appendix A of the user guide.
% Inserted before the end of the document by tools/build_guide_pdf.py, which also replaces the
% line that includes ported_functions.tex by the tables generated from docs/ported_functions.md.
\clearpage
\setcounter{secnumdepth}{3}
\appendix
\providecommand{\E}{\operatorname{E}}
\newcommand{\code}[1]{\texttt{\detokenize{#1}}}
\newcommand{\ind}{\mathbf{1}}
\newcommand{\bZ}{\mathbf{Z}}
\newcommand{\bX}{\mathbf{X}}
\newcommand{\diff}{\Delta}
\newcommand{\statasrc}[1]{\texttt{\detokenize{#1}}}

\section{Methodological note}\label{app:note}

This appendix is the methodological note of the package. It states what the two estimators
compute (Sections~\ref{note:model}--\ref{note:var}), how the package computes it and why it is
fast (Section~\ref{note:cumsum}), how the bootstrap inference is organised
(Section~\ref{note:boot}), where it departs from the algorithm of Gong and Seo
(Section~\ref{note:gs}), what \texttt{xthenreg} adds (Section~\ref{note:xthenreg}), which conventions of the
Stata commands are reproduced on purpose (Section~\ref{note:quirks}), what has been compared with Stata and how
(Section~\ref{note:valid}), and where every Stata function went (Section~\ref{note:map}).

Three kinds of statements are mixed and kept apart in the text. A \emph{model} statement comes from
the papers the commands implement (Seo and Shin 2016; Seo, Kim and Kim 2019; Gong and Seo 2026;
Nguyen and Lai 2026). A \emph{convention} is a behaviour of the Stata code that the package
reproduces so that the numbers agree; Section~\ref{note:quirks} lists them. An \emph{implementation}
statement says how the package obtains a result and is checked by tests, not by a theorem.
Nothing here claims more theory than the papers prove: where a scheme is an approximation that the
theory does not cover, the text says so.

\subsection{Notation}\label{note:notation}

\begin{longtable}{@{}p{2.6cm}p{13.4cm}@{}}
\toprule
Symbol & Meaning \\
\midrule
\endhead
$i=1,\dots,n$ & units (the clusters of the variance and of the bootstrap); $n$ is \code{e(N_units)} \\
$e$ & one transformed equation: a unit and a period. There are $N$ of them, \code{e(N)}, stacked by unit then period \\
$y_{it}, x_{it}, q_{it}$ & dependent variable, the $k$ regressors (the lag of $y$ first, in the dynamic model) and the threshold variable \\
$\gamma$ & the threshold; $\Gamma$ is the grid of candidates, $G$ its size \\
$\mathcal T$ & the transformation (first differences or forward orthogonal deviations) as a sparse matrix that maps the \emph{level} observations to the $N$ equations \\
$\theta=(\beta',\delta')'$ & coefficients; $p=2k+1$ in the jump model, $p=k+1$ in the kink model \\
$\bX(\gamma)$ & transformed regressors at $\gamma$: $\mathcal T\,[\,x,\ \ind(q>\gamma),\ x\,\ind(q>\gamma)\,]$ (jump) or $\mathcal T\,[\,x,\ (q-\gamma)\ind(q>\gamma)\,]$ (kink) \\
$\bZ$ & the $N\times L$ instrument matrix; $z_e$ is its row $e$ \\
$D(\gamma)$, $m_y$ & $D(\gamma)=\bZ'\bX(\gamma)/N$ ($L\times p$) and $m_y=\bZ'\,\mathcal T y/N$ ($L$) \\
$g_i(\theta,\gamma)$ & moment contribution of unit $i$: $\sum_{e\in i} z_e\,(\mathcal T y_e-\bX_e(\gamma)'\theta)$ \\
$W_1,\ W_2$ & weight matrices of the first and second step \\
$Q(\theta,\gamma)$ & criterion $N\,\bar g'\,W\,\bar g$ with $\bar g=m_y-D(\gamma)\theta$ \\
\bottomrule
\end{longtable}

Stata counts the stacked equations $N$ in \code{e(N)} and the units in \code{e(N_units)}; the note
follows that. (The papers write $n$ for the units and index the moments by period.)

\subsection{Model and moment conditions}\label{note:model}

The commands fit the single-threshold dynamic panel model of Seo and Shin (2016),
\begin{equation}\label{eq:jump}
y_{it}=x_{it}'\beta+(1,x_{it}')\,\delta\;\ind(q_{it}>\gamma)+\eta_i+\varepsilon_{it},
\end{equation}
where $x_{it}$ contains $y_{i,t-1}$ unless the model is static. The regression function jumps at
$\gamma$ by $\delta_1+x'\delta_2$. With \code{kink}, continuity is imposed,
\begin{equation}\label{eq:kink}
y_{it}=x_{it}'\beta+\delta_q\,(q_{it}-\gamma)\,\ind(q_{it}>\gamma)+\eta_i+\varepsilon_{it},
\end{equation}
and $q$ should be a regressor so that it has a slope below $\gamma$. The unit effect $\eta_i$ is
removed by a transformation.

\paragraph{First differences.} The equation of unit $i$ at period $t$ exists when the unit has rows
at the calendar periods $t$ and $t-1$ (a gap in the dates gives no equation) and both are in the
equation sample and complete in $y$, $q$ and the regressors. It is $y_t-y_{t-1}$ on
$x_t-x_{t-1}$, and touches two level observations with weights $+1$ and $-1$.

\paragraph{Forward orthogonal deviations.} An equation row followed by $T_f\ge1$ later equation rows of
the same unit is $c\,(y_t-\overline{y}_{\text{later}})$ with $c=\sqrt{T_f/(T_f+1)}$ and the same for
every regressor; the dates of the later rows do not matter, and the last row of a unit has no
equation. In weights: $+c$ on the row, $-c/T_f$ on each later row.

\paragraph{Triplets.} Both transformations are stored as triplets (equation, level observation,
weight), i.e.\ as a sparse matrix $\mathcal T$ with two entries per row for first differences. The
regime regressors depend on $\gamma$ through the level observations only, so $\bX(\gamma)$ is the
transform of the level regressors $[x,\ \ind(q>\gamma),\ x\,\ind(q>\gamma)]$: the regressors at any
threshold come from one sparse product, and the equations are never restacked. This is what makes
the cumulative-sum algorithm of Section~\ref{note:cumsum} possible.

\paragraph{Instruments.} The vector $z_e$ of an equation at period $t$ contains (Arellano and Bond 1991;
Roodman 2009), each in its own block of columns per period unless \code{collapse}:
a constant; under first differences the levels $y_{t-2},y_{t-3},\dots$; the transformed strictly
exogenous regressors; the lagged levels of the endogenous regressors from $t-2$ and of the
predetermined regressors from $t-1$; and the external instruments dated $t$. Under forward
orthogonal deviations the dates are one period later, since the transformed error involves later
periods. \code{maxlag(lo hi)} bounds the lags, the minimum being 2 for $y$ and the endogenous
variables and 1 for the predetermined ones. With \code{collapse} the per-period blocks are replaced
by one shared block (one column per lag instead of one per lag and period). An equation without an
instrument other than the constant is dropped. The instrument matrix does not depend on $\gamma$.

Columns that carry no information are then removed, in the order of the Stata command:
(1)~columns that are multiples of the period constants (a variable common to all units in a period)
are zeroed; (2)~columns that are exactly zero are dropped; (3)~exact duplicates are dropped, the
first kept; (4)~columns that are linear combinations of the earlier ones are dropped (squared
relative residual below $10^{-13}$). The counts are \code{e(N_iv_common)}, \code{e(N_iv_dep)} and
\code{e(N_iv_dep_near)}.

\paragraph{Moment conditions.} With $T(\cdot)$ the chosen transformation, the estimator uses
$\E[z_e\,\mathcal T(\varepsilon)_e]=0$. The sample moment is
$\bar g(\theta,\gamma)=m_y-D(\gamma)\theta$ and the criterion is
$Q(\theta,\gamma)=N\,\bar g'W\bar g$.

\paragraph{Time effects.} With \code{td} the transformed period dummies are partialled out of the
transformed $y$, regressors and instruments (Frisch--Waugh--Lovell). Under first differences the
transformed dummy of a period is the same for every unit, so the projection is exact demeaning within
the period; a period with one equation is demeaned to zero and its row is dropped. Under forward
orthogonal deviations the dummies receive the same transform as the data and the system is projected
on them with a generalised inverse (\code{invsym}); a row that the projection saturates is dropped.
The projections $M$ are symmetric and idempotent, so $(M\bZ)'\,d\bX=\bZ'\,M\,d\bX$ and the regime
moments are formed from the projected instruments and the unprojected level regressors.

\subsection{Estimation}\label{note:est}

\paragraph{Closed form at a threshold.} For a weight matrix $W$ and a given $\gamma$ the criterion is
quadratic in $\theta$, so
\begin{equation}\label{eq:theta}
\hat\theta(\gamma)=\bigl(D'WD\bigr)^{-1}D'W\,m_y,\qquad
\bar g=m_y-D\hat\theta,\qquad Q=N\,\bar g'W\bar g .
\end{equation}
The inverse is Mata's \code{invsym}, or the gated \code{syminv} of \code{xtdpthresh}
(Section~\ref{note:quirks}): a singular normal matrix is not an error but a candidate that is not
admissible.

\paragraph{First step.} $\theta_1$ and $\gamma_1$ minimise $Q$ over the grid with $W_1$. Under first
differences $W_1$ is the MA(1) weight of Arellano and Bond (1991): the transformed errors of
consecutive equations of a unit overlap in one term, so with $S_0=\bZ'\bZ/n$ and
$S_1=n^{-1}\sum_i\sum_t z_{i,t-1}z_{it}'$ over consecutive equations of the same unit,
\begin{equation}\label{eq:w1}
W_1=\bigl(2S_0-S_1-S_1'\bigr)^{-1}.
\end{equation}
If that matrix cannot be inverted, $(\bZ'\bZ)^{-1}$ is used, then the identity
(\code{e(W1_fallback)} $=1$ or $2$). Under forward orthogonal deviations $W_1=(\bZ'\bZ/N)^{-1}$.

\paragraph{Second step.} With the first-step residuals $r_1$ at $(\theta_1,\gamma_1)$, the cluster
covariance of the moments is
\begin{equation}\label{eq:omega}
\Omega=\frac1N\sum_{i=1}^n g_ig_i'-\frac{1}{nN}\Bigl(\sum_i g_i\Bigr)\Bigl(\sum_i g_i\Bigr)',
\end{equation}
the second term being dropped under \code{nocenter}. $W_2=\Omega^{-1}$ is built \emph{once}, from the
first-step estimate, and then \emph{held fixed} while the second-step criterion is evaluated at every
other threshold; $\hat\gamma$ minimises that criterion over the grid, with ties going to the smaller
threshold (relative tolerance $10^{-12}$). This is what the Stata search does, and it is why the
criterion of every candidate can be computed from sufficient statistics.

\paragraph{One-step fallback.} If $\Omega$ cannot be inverted, no candidate can be solved with $W_2$,
or the criterion does not vary over the admitted candidates of the second step, the first-step
estimates are reported with the sandwich of $W_1$ (\code{e(estimator_twostep)} $=0$) and without
Hansen's $J$. If the first-step criterion itself is flat, the command stops with return code 498.

\paragraph{The grid.} The effective support of $q$ is the set of $q$ values of the level observations
that enter a stacked equation (periods $t$ and $t-1$ under first differences); \code{gridsample(observed)}
keeps the current-row values only. The trimming bounds are the type-7 quantiles at
\code{trim}/2 and $1-$\code{trim}/2. A uniform grid has equally spaced points between the bounds;
a quantile grid places them on the empirical quantiles (duplicates collapsed). A candidate is
admitted when both regimes have at least $m$ observations, where the default floor is
$m=\max\{2,\min(\#\{q\le q_{lo}\},\#\{q>q_{hi}\})\}$ and \code{minregime} can only raise it, and when
there are more instruments than parameters. At least 20 equations, 5 units and 20 usable values of
$q$ are required. The grids of 100, 199 and 397 uniform points are nested ($g\to2g-1\to4g-3$), so
a finer grid is a direct check.

\paragraph{Refinement.} Under \code{kink} the criterion is continuous in $\gamma$, and $\hat\gamma$
minimises it between the two grid points next to the grid minimum: up to 8 rounds of 40 and then 10
trial points, shrinking the bracket until its spacing is below $10^{-6}$ of its width
(\code{e(kink_refined)}). Under the jump model \texttt{refine(\#)} adds the observed values of $q$
between the neighbours of the minimiser, in at most \# rounds, and re-minimises the second-step
criterion with the same fixed $W_2$; the confidence set and the tests keep using the initial grid.

\subsection{Cumulative sums: the criterion on a whole grid in one pass}\label{note:cumsum}

Only the regime columns of $\bX(\gamma)$ depend on $\gamma$. For instrument $j$ and level regressor
$k$ the regime moment is
\begin{equation}\label{eq:S}
S_{jk}(\gamma)=\sum_{e}\bZ_{ej}\sum_{(l,w)\in e} w\,W_{lk}\,\ind(q_l>\gamma),
\end{equation}
where the inner sum runs over the level observations $l$ of equation $e$ with their weights $w$
(for the jump model $W_l=(1,x_l)$). Sort the grid, $\gamma_1<\dots<\gamma_G$, and let
$m_l=\#\{g:\gamma_g<q_l\}$. Then $\ind(q_l>\gamma_g)=1$ exactly when $g\le m_l$, so each term
$c=\bZ_{ej}\,w\,W_{lk}$ contributes to the first $m_l$ grid points and to no other. The kernel adds
each contribution \emph{once}, to the bucket $m_l$, and a reverse cumulative sum over the buckets
gives every grid point:
\begin{equation}\label{eq:cumsum}
A[m]=\sum_{\{p:\,m_p=m\}}c_p,\qquad S(\gamma_g)=\sum_{m\ge g}A[m].
\end{equation}
The cost is $O(P\cdot\mathrm{nnz}\cdot K+G\,L\,K)$, with $P$ the number of pairs and $\mathrm{nnz}$
the number of non-zero entries of an instrument column, against $O(G\,P\,\mathrm{nnz}\,K)$ if each
threshold were recomputed. The kernel runs over instruments (one Numba thread owns the slice
$S[:,j,:]$), reads $\bZ$ in compressed-column form so that the block-sparse layout costs nothing, and
places a missing $q$ in bucket 0, so that it is never above the threshold.

For the kink model the regime column is $(q-\gamma)\ind(q>\gamma)=q\,\ind(q>\gamma)-\gamma\,\ind(q>\gamma)$,
so the same pass with level columns $(1,q)$ gives $S_\kappa(\gamma)=S_q(\gamma)-\gamma\,S_1(\gamma)$.
The full matrix is $D(\gamma_g)=[\,\bZ'\bX_{\text{lin}}/N\ \big|\ S(\gamma_g)/N\,]$, the first block
being the same at every threshold.

\paragraph{Solving the whole grid.} Given $D(\gamma_g)$ for all $g$, the estimates (\ref{eq:theta}) are
batched matrix products (BLAS) and one batched \code{invsym}; this step is dominated by dense algebra
and is plain NumPy. On the $565\times15$ panel the regime moments of 100 grid points take 4~ms
against 278~ms for the straightforward reference implementation that is kept for the equivalence
tests (block~C of the project plan), and the two agree to $10^{-12}$. At realistic sizes the
remaining time is in the bootstrap (Section~\ref{note:boot}), whose criterion is one matrix product
that already runs close to the speed of the machine, which is why no further algebraic gain exists
there: an idea of sharing the unrestricted fit across candidates was examined and rejected because the
residual depends on the candidate.

\paragraph{Unit moments.} $g_i=\sum_{e\in i}z_e r_e$ is a segmented sum over the CSR layout of $\bZ$
(one kernel). Units without equations are zero rows and change nothing in (\ref{eq:omega}).

\subsection{Variances and specification tests}\label{note:var}

\paragraph{Robust sandwich.} For a fixed weight $A$ (efficient or not), with $D=\bZ'\bX/N$ and
$B=(D'AD)^{-1}$,
\begin{equation}\label{eq:sandwich}
\widehat V=\frac1N\,B\,D'A\,\Omega\,A D\,B,
\end{equation}
with $\Omega$ rebuilt from the residuals at the reported estimate. It is valid for the one-step
estimator too. This is the variance \emph{conditional on $\hat\gamma$}.

\paragraph{Joint variance with the threshold.} Following Seo and Shin (2016) the threshold is one more
parameter. The Jacobian is augmented to $[\,D\ |\ \bZ'x_\gamma/N\,]$, where $x_\gamma$ is the
derivative of the transformed residual with respect to $\gamma$, and the $\theta$ block of the
sandwich of the augmented system is reported. In the jump model the sample moment is a step function
of $\gamma$, so the derivative of its expectation is smoothed by a Gaussian kernel,
\begin{equation}\label{eq:xg}
x_\gamma=-\,\mathcal T\Bigl[\varphi_h(q-\hat\gamma)\,(\hat\delta_1+x'\hat\delta_2)\Bigr],\qquad
h=c\cdot1.06\cdot s_q\cdot n^{-1/5},
\end{equation}
with $\varphi_h(u)=\varphi(u/h)/h$, $s_q$ the standard deviation of $q$ over all level observations,
$n$ the number of units with an equation and $c=$ \code{bwscale} (1.5 by default, as in
\code{xthenreg}). In the kink model the derivative is analytic, $x_\gamma=-\hat\delta_q\,\mathcal T\,\ind(q>\hat\gamma)$.
The variance conditional on $\hat\gamma$ is kept (\code{e(V_cond)}); if the joint variance cannot be
computed, the conditional one is reported (\code{e(joint_vce)} $=0$). The upper-regime slopes are
$\hat\beta+\hat\delta_2$, with the variance of the linear combination.

\paragraph{Windmeijer correction.} $W_2=\Omega(\theta_1)^{-1}$ depends on the first-step estimate.
Linearising the second-step estimate in $\theta_1$ (Windmeijer 2005) adds
$D_j=-B_2G_2'W_2\,[\partial\Omega/\partial\theta_j]\,W_2\bar g_2$, with $B_2=(G_2'W_2G_2)^{-1}$,
$G_2=D(\hat\gamma)$, $\bar g_2$ the second-step moment residual and $B_1,G_1$ the first-step
counterparts. The variance is
\begin{equation}\label{eq:wind}
\widehat V_W=\frac1N\,M\,\Omega_1\,M',\qquad M=B_2G_2'W_2+[D_1,\dots,D_p]\;B_1G_1'W_1 .
\end{equation}
The derivative of $\Omega$ is never formed. With $h_{ij}=\sum_{e\in i}z_ex_{ej}$ and
$s=\sum_ig_i$, only the product
$[\partial\Omega/\partial\theta_j]\,v=-\frac1N\sum_i\bigl(h_{ij}\,g_i'v+g_i\,h_{ij}'v\bigr)
+\frac{1}{nN}\bigl(h_j\,s'v+s\,h_j'v\bigr)$ with $v=W_2\bar g_2$ is needed (the last term comes from the
centring), which costs two products per column instead of an $L\times L$ matrix. The correction is
computed once at the final estimate; if it cannot be computed or is not positive semidefinite, the
robust variance is reported (\code{e(vce_applied)} $=0$). Windmeijer derived it for moments linear
in the parameters; it need not improve the variance for the nonlinear moments of the threshold.

\paragraph{Hansen's $J$.} $J=Q(\hat\theta,\hat\gamma)$ of the second step, referred to $\chi^2$ with
$L-p-1$ degrees of freedom (the threshold counts as one parameter). Nothing is reported after the
one-step fallback or when the model is not over-identified.

\paragraph{Arellano--Bond serial correlation.} On the \emph{first-difference} residuals, whatever the
estimation transformation (the \code{xtabond2} convention), the statistic of order $m$ is
$\mathrm{AR}(m)=b_0/\sqrt{T_1+T_2+T_3}$ (Arellano and Bond 1991, eq.~8), where, with
$r_{it}$ the residuals and $u_i=\sum_t r_{it}r_{i,t-m}$,
$b_0=\sum_iu_i$, $T_1=\sum_iu_i^2$, $T_2=-2\,g'(MAM')^{-1}MA\,s$ and $T_3=g'\widehat V g$. Here
$g=\sum_e x^*_e\,r_{e,t-m}$ collects the regressors against the lagged residual, $s=\bZ'(r\circ u_{\text{unit}})$
the instrument moments weighted by the unit sums, $M=\bZ'\bX$ (raw sums) and $\widehat V$ the variance
of the parameters. With the joint variance the threshold column $x_\gamma$ is a parameter of $g$, $M$ and
$\widehat V$ (\code{e(ar_joint)} $=1$); the statistics that treat $\hat\gamma$ as known are
\code{e(ar1_cond)} and \code{e(ar2_cond)}. The p-value is $2\Phi(-|\mathrm{AR}|)$. At least five
residual pairs, from at least five units, are needed.

\subsection{Bootstrap inference}\label{note:boot}

\paragraph{The wild bootstrap of the command.} Every bootstrap component compares, on each sample, the
criterion of a \emph{restricted} model with the smallest criterion of the threshold models over the
initial grid. A sample is $Y^*_b=F+r\,\eta_{i(e),b}$: a fitted part $F$ and the residuals $r$ of the
restricted fit, multiplied by one weight $\eta_{ib}$ per unit (so the correlation within a unit and
the heteroskedasticity across units are kept), the same weights for every candidate. Regressors and
instruments keep their sample values. The weights follow Mammen (1993): $-\phi$ with probability
$(\sqrt5+1)/(2\sqrt5)$ and $1/\phi$ otherwise, $\phi=(\sqrt5-1)/2$; mean 0, variance 1.

\paragraph{Why it is cheap.} With a fixed weight $W$ the criterion of a candidate $g$ depends on the
sample only through the moments $m^*=\bZ'Y^*/N$:
\begin{equation}\label{eq:quad}
\theta=C_g\,m^*,\quad C_g=(D_g'WD_g)^{-1}D_g'W,\qquad
J_g(m^*)=N\,{m^*}'\,\bigl[M_g'WM_g\bigr]\,m^*,\quad M_g=I-D_gC_g .
\end{equation}
The moments of all samples are $m^*_b=\bZ'F/N+N^{-1}\sum_i\eta_{ib}\,g_i$, where $g_i=\sum_{e\in i}z_er_e$
is the unit moment of the residuals. No $N\times B$ matrix is ever formed: for $G$ candidates and $B$
draws, all the criteria are a few batched products of the $G$ matrices
$Q_g=M_g'WM_g$ ($L\times L$) with the $L\times B$ matrix of moments, in chunks that bound the memory
(8 million elements for the work array; the $Q_g$ are kept when they fit in 40 million elements and
the estimates are solved draw by draw otherwise). The bootstrap is 80--92~\% of the time of a full
run; it is evaluated this way for all draws and candidates at once.

\paragraph{Confidence set (grid bootstrap).} Following Hansen (1999b) and Gong and Seo (2026, \S4.1), the
set collects the candidates $\gamma_\ell$ that a bootstrap test does not reject. With the weight of
the reported estimator ($W_2$ after a two-step fit, $W_1$ after the fallback), for each candidate:
(1)~restricted fit $\hat\theta(\gamma_\ell)$ and residuals $e$; (2)~statistic
$D_n(\gamma_\ell)=Q(\hat\theta(\gamma_\ell),\gamma_\ell)-m_n$, where $m_n$ is the minimum of
$Q(\hat\theta(\gamma),\gamma)$ over $\Gamma\cup\{\gamma_\ell\}$ (so $D_n(\hat\gamma)=0$);
(3)~draws $Y^*_b=\bX(\gamma_\ell)\hat\theta(\gamma_\ell)+\eta_{ib}e$; (4)~$D^*_{n,b}(\gamma_\ell)=J_\ell(m^*_b)-\min\{J_\ell(m^*_b),\min_{g\in\Gamma}J_g(m^*_b)\}$;
(5)~accept $\gamma_\ell$ iff $D_n(\gamma_\ell)$ does not exceed the order statistic of rank
$k=\lceil p\,(B+1)\rceil$ of the $D^*_{n,b}$, with $p=$ \code{level}/100. That is the add-one rule:
the bootstrap p-value $(1+R)/(1+B)$, $R=\#\{D^*\ge D\}$, exceeds $1-p$. If $k>B$ no draw can reject and
every candidate is accepted. The set is the union of the accepted candidates; \code{e(ci_nseg)}
counts its segments and \code{e(gamma_lo)}, \code{e(gamma_hi)} are the hull, which contains rejected
values when there are several segments. A candidate whose restricted fit cannot be solved is not
admissible (status 3); a draw that gives no valid statistic leaves the candidate unresolved
(statuses 4--6), and then no set is reported (\code{e(ci_incomplete)} $=1$). A statistic exactly zero
is accepted mechanically (status 2).

\paragraph{Linearity test.} $H_0:\delta=0$. The statistic is the one-step criterion of the linear model
minus the smallest one-step criterion of the threshold model over the initial grid, both with $W_1$;
samples are drawn from the linear fit; the p-value is $(1+R)/(1+B)$, and is missing if any draw is invalid.
This GMM distance test differs from the sup-Wald test of \code{xthenreg}.

\paragraph{Test of a given threshold.} \texttt{citest(\#)} is one row of the grid bootstrap at $\gamma=\#$,
with its own seed, so asking for it changes no other result; its p-value is $(1+R)/(1+B)$ and its
decision is $D\le$ crit.

\paragraph{Continuity test.} Following Gong and Seo (2026, \S3.2 and \S4.3) with
$\tilde\theta=\arg\min_{\Theta_c}\hat Q_n$ the continuity-restricted estimator, the statistic is the
distance $T_n=N(Q(\tilde\theta)-Q(\hat\theta))$ between the kink and the jump criteria. The package computes
it over the initial grid with the common weight, draws samples from the kink fit with the \emph{jump}
residuals reweighted by unit, and takes the p-value $(1+R)/(1+B)$ of $T^*_n=\min_{\text{kink}}J^*-\min_{\text{jump}}J^*$.
It needs $q$ as a contemporaneous regressor, otherwise the kink model is not nested in the jump model.

\paragraph{Coefficient bootstrap (legacy).} \code{coef_boot} re-runs the whole search in each draw: samples
$Y^*=\bX(\hat\gamma)\hat\theta+r\eta$, the threshold searched again over the estimation grid with $W_1$, $\theta$
re-estimated there (two-step: $W_2^*$ from the draw's residuals, threshold searched again). The intervals are
$\hat\theta_j\pm q_{1-\alpha}(|\theta^*_j-\hat\theta_j|)$ (symmetric) or the percentiles
$q_{\alpha/2},q_{1-\alpha/2}$ of $\theta^*_j$. Section~\ref{note:gs} says why this is not the coefficient
bootstrap of Gong and Seo.

\paragraph{Unit resampling.} \code{boot_type="unit"} follows the structure of Algorithm 1
(Section~\ref{note:gs}). A resample is a vector of unit counts, so every quantity is a sum over units of
per-unit cross products that do not depend on the draw: $\bZ'\bX(\gamma_g)$ of each grid point, $\bZ'y$ and the
unit moments; a draw costs a few small products and no pass over the $N$ equations. Per draw: the
first-step coefficients of every grid point with $W_1$ of the sample, the best one giving the residuals whose
recentred cluster covariance defines $W_2^*=\Omega^{*-1}$; then the criterion at the candidate (restricted)
and at every grid point (unrestricted) with $W_2^*$. A draw whose profile has no usable minimum is skipped,
not replaced.

\paragraph{Random numbers.} Stata 17 uses the 64-bit Mersenne Twister; \code{set seed s} starts it with
\code{init_genrand64(s)}, and a draw is $u=((x\gg11)\,|\,1)\cdot2^{-53}$, where $x$ is the next 64-bit output
(measured on 10\,000 draws from five seeds: identical). A matrix is filled row by row;
\code{uniform()}, \code{runiform()} and \code{invnormal(uniform())} consume one stream. The package reproduces
the stream, so the same seed gives the same weights as the command. With \code{rseed(s)} each component restarts
at $(s+\text{offset})\bmod2^{31}$: confidence set $0$, linearity $104\,729$, continuity $224\,737$, coefficient
$350\,377$, \code{citest} $477\,377$. The weights of a component are drawn replication after replication,
one uniform per contributing unit. Without a seed one random seed is drawn and the components share a
stream in the order they run. The Python wrapper always records a seed (Section~\ref{note:quirks}).

\subsection{Gong and Seo's Algorithm 1 and what the package does}\label{note:gs}

This section is the reading of Gong and Seo (2026), Algorithm~1 and \S\S4.1--4.3 (Theorems 5--7), against the code.
The proofs were not re-derived; the comparison is of the algorithms, and the statements about validity are the
paper's.

\paragraph{The algorithm of the paper.} For a parameter $\theta^*_0=(\beta^{*\prime}_0,\delta^{*\prime}_0,\gamma^*_0)'$ that
generates the bootstrap data:
\begin{enumerate}
\item draw $i^*$ i.i.d.\ uniformly from $\{1,\dots,n\}$ and resample the units \emph{jointly}:
$(x^*_{it},x^*_{it-1},z^*_{it},\widehat{\Delta\epsilon}^*_{it})=(x_{i^*t},x_{i^*t-1},z_{i^*t},\widehat{\Delta\epsilon}_{i^*t})$,
with the \emph{unrestricted} residual $\widehat{\Delta\epsilon}_{it}=\Delta y_{it}-\Delta x_{it}'\hat\beta-\ind_{it}(\hat\gamma)'X_{it}\hat\delta$;
\item generate $\Delta y^*_{it}=\Delta x^{*\prime}_{it}\beta^*_0+\ind^*_{it}(\gamma^*_0)'X^*_{it}\delta^*_0+\widehat{\Delta\epsilon}^*_{it}$;
\item $g^*_{it}(\theta)=z^*_{it}(\Delta y^*_{it}-\Delta x^{*\prime}_{it}\beta-\ind^*_{it}(\gamma)'X^*_{it}\delta)$;
\item recentre: $\bar g^*_n(\theta)=\frac1n\sum_i\bigl(g^*_i(\theta)-\bar g_n(\hat\theta)\bigr)$, with
$\bar g_n(\hat\theta)$ the sample moment of the unrestricted residuals;
\item $\hat\theta^*_{(1)}=\arg\min\bar g^{*\prime}_n\bar g^*_n$ (identity weight) and
$W^*_n=\bigl(\frac1n\sum g^*_ig^{*\prime}_i-\bar g^*\bar g^{*\prime}\bigr)^{-1}$ at $\hat\theta^*_{(1)}$;
\item $\hat Q^*_n(\theta)=\bar g^*_n(\theta)'W^*_n\bar g^*_n(\theta)$ and the bootstrap statistic.
\end{enumerate}
Three choices of $\theta^*_0$ give three procedures. The grid bootstrap of the threshold takes
$\theta^*_0=(\hat\alpha(\gamma_\ell)',\gamma_\ell)'$ (the restricted estimate at the candidate) and
$D^*_n(\gamma_\ell)=n\bigl(\min_\alpha\hat Q^*_n(\alpha,\gamma_\ell)-\min_\theta\hat Q^*_n(\theta)\bigr)$;
the set is $\{\gamma\in\Gamma_n:D_n(\gamma)\le\hat F^{*-1}(1-\tau;D^*_n(\gamma))\}$ (eq.~7), which is never empty because
$D_n(\hat\gamma)=0$. Theorem~5: at $\gamma=\gamma_0$, $D^*_n\to Z_0^2$ with $Z_0=\max(0,Z^*_0)$ if the model is continuous and
$\to\chi^2_1$ if it is not, so the set is valid whether or not the function is continuous. The coefficient
bootstrap takes $\theta^*_0=w_n\hat\theta+(1-w_n)\tilde\theta$ with $w_n=\min\{T_n/(\hat Cn^{1/4}),1\}$ and the
continuity-restricted estimator $\tilde\theta$ (eq.~8), so that the bootstrap Jacobian degenerates fast enough
under continuity; the intervals are eq.~(9) (percentile) or eq.~(10) (symmetric), Theorem~6, pointwise only. The continuity test
takes $\theta^*_0=\tilde\theta$ and $T^*_n=n(\min_{\Theta_c}\hat Q^*_n-\min_\Theta\hat Q^*_n)$ (Theorem~7).

\paragraph{Comparison.}
\begin{longtable}{@{}p{3.1cm}p{6.1cm}p{6.8cm}@{}}
\toprule
Step & Gong and Seo, Algorithm 1 & \code{boot_type="wild"} (default) \\
\midrule
\endhead
Resampling & units drawn with replacement; $x,z$ and the residual move together & one Mammen weight per unit multiplies the residuals; $x$ and $z$ keep their sample values \\
Data-generating process & $\theta^*_0$ and \emph{unrestricted} residuals & restricted fit at the candidate and \emph{restricted} residuals; $\E^*[\eta]=0$, so the null holds without recentring \\
Recentring & subtract $\bar g_n(\hat\theta)$ & not needed (restricted residuals, mean-zero weights) \\
Weight in the bootstrap criterion & re-estimated in every draw from a first step with the identity & the sample weight, held fixed ($W_2$, or $W_1$ after the fallback) \\
Validity & Theorem 5 (continuous or not) & not covered by the theorem; coverage is examined by simulation (Nguyen and Lai 2026) \\
\bottomrule
\end{longtable}

\noindent \code{boot_type="unit"} follows the structure instead: whole units resampled (as count vectors),
restricted coefficients $\hat\theta(\gamma_\ell)$ with the \emph{unrestricted} residuals at $(\hat\theta,\hat\gamma)$, moments recentred at the sample moment
of those residuals, and $W^*_2$ re-estimated in every draw. One difference is deliberate in the Stata command and
copied: the first step of each draw uses the sample first-step weight $W_1$ where the paper uses the identity. It
is available for \code{method="fd"} without \code{kink} or \code{td}, and is much slower. An unbalanced effective
stack takes the command outside the common-$T$ setting of the theory (Section~\ref{note:map} lists the warning
that Stata prints and the package does not).

\paragraph{The coefficient bootstrap is not the paper's.} The legacy \code{coef_boot} generates data from
$\bX(\hat\gamma)\hat\theta$ with reweighted residuals, i.e.\ $\theta^*_0=\hat\theta$, which is the choice $w_n=1$ of
eq.~(8). That is the right choice when the model is discontinuous; under continuity the paper shows that a
bootstrap built on $\hat\theta$ is not valid, which is why it proposes $w_n$. The command therefore calls it a
legacy compatibility feature, \code{coef_boot="gs"} is refused, and the package says so in a note when it is used.
The symmetric interval is eq.~(10) with $\theta^*_0=\hat\theta$.

\paragraph{What is open.} The comparison above is a reading, not a proof of agreement. The power of the
rejection of shifted thresholds in the package's Monte Carlo experiments on the designs of Gong and Seo is lower than
their Table~2 while the coverage of the threshold agrees, and with the self-exciting designs of Seo and Shin
(2016) the dispersion of the estimated threshold is larger than in their tables, although the extreme estimates are
identical to Stata's, so they come from the estimator. The cause is not established.

\subsection{\texttt{xthenreg}}\label{note:xthenreg}

\paragraph{Estimator.} The balanced-panel estimator of Seo, Kim and Kim (2019) is the first-difference special
case of Section~\ref{note:est}: equation periods $t=2,\dots,T$, a block of instruments per period
(constant, $\diff x_t$ of the plain regressors, the levels $y_1,\dots,y_{t-2}$, the levels of the endogenous
regressors up to $t-2$, the external instruments), the coefficient order $[\beta\,|\,\delta_1\,|\,\delta]$ for the
jump and $[\beta\,|\,\kappa]$ for the kink, the first-step MA(1) weight built block by block
($W_1=(2S_0-S_1-S_1')^{-1}$ with $S_0$ block diagonal, $(1/N)Z_p'Z_p$, and $S_1$ block upper diagonal, $(1/N)Z_{p-1}'Z_p$), and the
second-step weight from the centred covariance of the unit moments. The grid is uniform between two
\code{mm_quantile} bounds at \code{trim_rate}/2 and $1-$\code{trim_rate}/2 over all the values of $q$ (default
\code{trim_rate} $0.4$) and $\hat\gamma$ is the first minimum of the second-step criterion; both steps
search the grid, the second step with the weight built at the first-step estimate.

\paragraph{Covariance.} Following Seo and Shin (2016, p.~173), with $G=(G_b,G_d,G_r)$ the derivatives of the
moments with respect to the linear coefficients, the regime coefficients and the threshold,
\begin{equation}\label{eq:xhcov}
\widehat V=\frac1N\,\bigl(G'\,\Omega^{-1}\,G\bigr)^{-1},
\end{equation}
with the derivative in $\gamma$ computed by the same Gaussian kernel as in (\ref{eq:xg}), bandwidth
$h=h_0\cdot1.06\cdot N^{-1/5}\cdot\mathrm{sd}(q)$ (\code{h_0(1.5)}), in the jump model, and in closed form for the kink.
The 95\,\% interval of every coefficient, the threshold included, is the estimate $\pm1.959963984540054\,\mathrm{se}$.

\paragraph{Linearity test.} The Wald statistic of the regime coefficients is computed at every grid point, each time
with its own two-step estimate and covariance, and the statistic is the supremum over the grid
(\code{e(sup_wald)} in the package). The p-value is from the fast bootstrap of the paper: the step-two moments of each
unit are multiplied by a standard normal weight ($\texttt{invnormal(uniform())}$), the one-step estimator is
re-evaluated with the fixed step-two weight at every grid point, and the Wald form reuses the covariance matrices of the
sample. The p-value is the \emph{share} of bootstrap suprema above the sample supremum, not the add-one rule of
\code{xtdpthresh}, so it can be exactly zero; a missing bootstrap statistic counts as exceeding.

\paragraph{Against \texttt{xtdpthresh}.} Under \code{method(fd)} the two commands use the same moment conditions and
weight matrices and, on the same grid, give the same estimates; the differences are in the trimming default, the
placement of $q$, the exogenous-regressor instruments (levels in \code{xthenreg}, transformed values in
\code{xtdpthresh}), the refinement of the kink, the absence of a standard error for $\hat\gamma$ and the GMM-distance
linearity test of \code{xtdpthresh}. Where the two designs coincide (cases c1, c4 and c5 of the validation) the package
checks that the two ports agree.

\subsection{Conventions of the Stata commands that are reproduced}\label{note:quirks}

Each row is a behaviour of the Stata code that the package copies so that the numbers agree. The ones marked
\code{STATA-QUIRK} in the source are the ones where the behaviour is surprising.

\begin{longtable}{@{}p{3.6cm}p{12.4cm}@{}}
\toprule
Convention & What is reproduced and why \\
\midrule
\endhead
\code{invsym} &
Reads only the upper triangle of its argument; a sweep with the pivot on the largest remaining diagonal element; a pivot
is dropped when it is at most $10^{-9}$ times the \emph{original} diagonal element, at most $9\times10^{-20}$, or negative,
and its row and column of the result are zero (a generalised inverse). Measured on Stata 17 MP. For a well-conditioned
positive definite matrix no pivot can be dropped in any order, so a Cholesky factorisation replaces the sweep. \\
\code{cholinv} &
Reads only the lower triangle; missing if the factorisation fails or the reciprocal condition number is below $10^{-20}$. \\
\code{syminv} gate &
\code{xtdpthresh} admits an inverse only after Jacobi equilibration, a Cholesky factorisation, a condition number of at most
$10^{12}$ and a check that $CC^{-1}\approx I$. A matrix that fails is a candidate that is not admissible. Validated on
23 matrices against the Mata source run unchanged. \\
Tolerances of ties &
Equal criteria (relative $10^{-12}$ for profiles, $10^{-10}$ for the refinement check) go to the smaller threshold
(\code{objtol}). \\
$W_2$ built once &
The second-step weight is built at the first-step threshold and held fixed over the grid of the second step. \\
Equal splits, exact zero &
Two candidates with the same split of the observed values of $q$ give the same regressors; the statistic of the estimated
threshold must be exactly zero, so the moments of the initial-grid entry with the same split are reused instead of
summing the arrays in another order (a gap of the order of rounding would otherwise remain). \\
Order statistic &
The critical value is the order statistic of rank $\lceil p(B+1)\rceil$, which is the add-one p-value rule; infinite when the
rank exceeds $B$ (no draw can reject). \\
Random stream &
Mersenne Twister 64, $u=((x\gg11)|1)2^{-53}$, row-major fill, component seeds $(s+\text{offset})\bmod2^{31}$, one draw per
contributing unit (a unit without equation takes no draw, so adding such a unit leaves the draws unchanged). \\
Unit bootstrap draws &
A draw whose profile has no usable minimum is skipped, not replaced (\code{STATA-QUIRK}); the missing ones stay missing in
the result. \\
\code{dlt_len} in \code{xthenreg} &
With a kink the Wald block that is tested has \emph{two} entries (the last linear coefficient and the kink
coefficient) instead of the kink coefficient alone (\code{STATA-QUIRK}); kept so that the statistic equals Stata's. \\
Signs of $G$ &
With a kink $G_d$ has the positive sign (unlike the jump model) and $G_r$ uses the indicator of the regime above the
threshold, as in \code{cov_mat_con_kink} (\code{STATA-QUIRK}). \\
\code{exogenous()} of \code{xthenreg} &
An \code{exogenous()} regressor enters with its levels up to $t$, two more lags than an \code{endogenous()} one
(\code{STATA-QUIRK}); hence 87 moment conditions where \code{xtdpthresh} has 60 instruments for the same data. \\
\code{boots_p} &
The share of bootstrap suprema above the sample supremum, a missing one counting as exceeding, as a missing value does in
Mata; $-1$ when the test is not run. \\
$h$ and $s_q$ &
$s_q$ is the sample standard deviation over \emph{all} the level observations and $n$ the units with an equation; the
1.959963984540054 of the interval is the constant typed in the \code{xthenreg} source. \\
AR on first differences &
The serial-correlation tests use first-difference residuals under both transformations. \\
Time-effect zeroing &
A column that the projection reduces to rounding noise (below $10^{-10}$ of its size) is set to exactly zero, as Stata does. \\
\code{if}, missing values &
A missing value in the condition counts as \emph{true} in Stata; in the package it counts as false, and the rows with a missing
variable are removed afterwards by the missing-value check, so the sample is the same. \\
Seed (the package adds) &
When the bootstrap is on and no seed is given, the wrapper draws one in $[1,2^{31}-1)$ and records it (\code{res.seed}, the
settings log, the JSON). Stata draws from the current generator and does not report it. \\
Precision &
All values are \code{float64}; Stata's operators on \code{float} variables are recomputed in double precision in the command
(\code{_xdpt_tsdouble}), which a double-precision array makes unnecessary. \\
\bottomrule
\end{longtable}

\subsection{Validation journal}\label{note:valid}

\paragraph{Rule.} Every number the package reports must equal Stata 17 MP to four decimals: absolute difference at most
$5\times10^{-5}$ \emph{or} relative difference at most $10^{-6}$. The comparison is automated and runs without Stata.

\paragraph{Method.} The real commands (\code{xthenreg}, \code{xtdpthresh} 0.9.38) are run in batch from PowerShell; a Mata
routine exports each object at 15 significant digits; the results are frozen in \code{tests/golden/} (one file per
component) with the versions of Stata and of the commands. The test suite compares the package with them case by case.
The random numbers of Stata are reproduced, so a seeded bootstrap draws the very weights of the command and its
p-values and sets are compared directly, not through their distribution. Each case is described once and drives both the
Stata run and the Python run.

\begin{longtable}{@{}p{4.3cm}p{2.2cm}p{9.5cm}@{}}
\toprule
Object & Cases & Result \\
\midrule
\endhead
Random numbers & 5 seeds & $10\,000$ draws each, identical; $40$ bootstrap suprema of each of 9 \code{xthenreg} cases equal to a relative $2\times10^{-10}$ \\
\code{invsym}, \code{cholinv}, \code{syminv} & measured; 23 matrices & thresholds measured on Stata; \code{syminv} equal to the Mata source run unchanged \\
\code{xthenreg}: $L$, $y$, $q$, $x$, $Z$ & 12 & equal to $10^{-12}$ running the Mata source as is (step~1 of \code{maincode_0719}) \\
\code{xthenreg}: $\hat\theta$, $\hat\gamma$, $V$, CI, $e()$ & c1--c8 and variants & $b$ within $3\times10^{-11}$ worst case on c1--c6; \code{h_0(3)} and \code{h_0(0.7)} (c7, c8) through the public function; grid $\times$ trimming variants (12 jump, 3 kink, 4 on $T=15$ with $L=325$) \\
\code{xthenreg}: sup-Wald, $p$ & 9; 5 & statistic to a relative $5\times10^{-10}$; the p-values of \code{boost(100)} equal on 5 cases (static and dynamic, jump and kink) \\
\code{xtdpthresh}: sample selection & 36 & \code{e(N)}, \code{e(N_stack)}, \code{e(N_units)} and the exact set of rows of \code{e(sample)} equal; return codes of the refusals equal \\
Instruments $Z$, $X$ & 28 + 30 + 15 & default, \code{maxlag}, \code{collapse}, \code{iv()}, duplicates, dependence, common variable; 15 cases of forward deviations \\
One step, two steps & 28 & \code{e(b)}, \code{e(obj)}, \code{e(V_cond)}; the one-step fallback is reproduced \\
Grid and search & 29; 5; 3 & \code{e(grid_admitted)}, \code{e(q_lo)}, \code{e(q_hi)}; kink refinement (\code{e(kink_refined)}); \code{refine} (iterations and points added) \\
Whole commands, full panel & 14 & the examples of the help file and variants on 565 firms: \code{e(b)}, \code{e(gamma)}, \code{e(obj)}, \code{e(N)}, \code{e(N_iv)} \\
Unbalanced panels, FOD, $td$, history & 5 + 5; 7 + 8; 4 & interior gaps, missing values, short units, alternate years; \code{td} under FD and FOD; \code{history(panel|sample)} \\
Variances & 45; 23 & \code{e(V)}, \code{e(V_cond)}, \code{e(V_upper)}, \code{e(b_upper)}, jump and kink, FD and FOD, \code{td}, unbalanced; \code{vce(windmeijer)}; \code{bwscale} (0.8, 3), \code{nocenter} \\
Tests & 45 & \code{e(hansen)}, \code{e(hansen_df)}, \code{e(hansen_p)}, \code{e(ar1)}, \code{e(ar2)} and their p-values \\
\code{predict} & 11 & residuals, \code{xb}, AR residuals, regime, \code{if} \\
Bootstrap & 62 Stata cases, 287 checks & \code{e(pval_lin)} (56), \code{e(ci_grid)}, \code{e(gamma_lo)}, \code{e(gamma_hi)}, \code{e(ci_nseg)}, \code{e(ci_segments)}; continuity (8); \code{citest} (8); coefficient bootstrap (10); \code{boottype(unit)}; refusals \\
Reporting & 146 checks & the tables repeat the validated $b$, $V$, \code{ci_grid}, \code{e(b_upper)} exactly \\
Hostile inputs & 19 & minimal panels, a constant, binary or rare threshold, an empty regime, duplicated regressors, fewer units than instruments, a null outcome: the same return code as Stata on 19 of 19\\
\bottomrule
\end{longtable}

\paragraph{Equivalence of the fast kernels.} Each Numba kernel has a NumPy reference that is kept; the tests compare them to
$10^{-12}$ (several seeds, ties between $q$ and the grid, a missing $q$, degenerate grids, one against four threads bit for
bit, and with the compiler switched off). The optimisations of the bootstrap were accepted only when the reference
cases remained identical.

\paragraph{Suite.} About 2\,170 tests (unit, golden, integration), line and branch coverage about 91.5\,\% (the compiled bodies
cannot be traced by the coverage tool); the code of the README and of this guide is executed by the tests.

\paragraph{What is not compared with Stata.} The statistical behaviour of the package's own seeds (size and power of the
linearity test, coverage of the confidence set over many panels) is measured by Monte Carlo and compared with the
papers, not with Stata, and has the open points of Section~\ref{note:gs}. Timings are measured on the same machine one run
after the other, because the speed of a machine varies by up to a factor of two from one moment to the next. The statements
of the guide about interpretation come from the Stata help files, the code and the execution of the package, and are not
a second validation.

\subsection{Where every Stata function went}\label{note:map}

The inventory is generated by \code{tools/ported_functions.py} from the \texttt{\# Stata: function (file.ado:line)} comments of
the source and the function ranges of \code{validation/stata_functions.json} (a manifest with the SHA-256 of the three
\code{.ado} files); a test checks that every citation points inside a function. A function counts as \emph{ported} when
at least one comment points inside it. At the time of writing, 86 of the 101 functions of the three files are cited; the generated
tables below list the cited ones and where they live. The 15 that are not cited fall in three groups.

\begin{longtable}{@{}p{3.7cm}p{3.0cm}p{9.3cm}@{}}
\toprule
Stata function & Status & Reason \\
\midrule
\endhead
\code{xdpt2_append_nested_level}, \code{xdpt2_nested_midpoints}, \code{xdpt2_profile_stability} &
retired &
Reachable only when \code{search_mode == "adaptive"}, which the command refuses since 0.9.26 (\code{search_mode}, \code{search_max} and
\code{search_tol} are retired). The package evaluates the grid it is asked for; the nested grids $g\to2g-1\to4g-3$ are reached by choosing
\code{grid=199} or \code{397}. \\
\code{xdpt2_build_gamma_cache} &
dead code &
No caller in the command; the variant \code{xdpt2_build_gamma_cache_t} is the one in use and is ported (the cache of the candidates). \\
\code{xdpt2_tpl_build}, \code{xdpt2_tpl_dW}, \code{xdpt2_tpl_release}, \code{xdpt2_unit_cfg} &
replaced &
The Mata record of the stacked-equation template, its release and the settings read once per stack. The package keeps the
transformation as the sparse matrix $\mathcal T$ (Section~\ref{note:model}) and the instrument settings as a dataclass, so none of this
bookkeeping exists; the quantities they produce are the ones compared with Stata. \\
\code{xdpt2_hasmiss} &
not needed &
Tests for a missing value in a block; the package uses \code{np.isfinite}. \\
\code{_xdpt_tsdouble} &
not needed &
Recomputes the time-series operators in double precision (\code{tsrevar} keeps the type of the source); all values are \code{float64}. \\
\code{xdpt2_rowsig2} (in \code{xtdpthresh.ado} and in \code{xtdpthresh_p.ado}), \code{xdpt2_p_fill} of \code{xtdpthresh_p.ado} &
replaced &
A checksum of the data and a keyed store that let \code{predict} detect a changed dataset or a cleared memory (return codes 459 and
498). The result object holds its data and \code{predict(kind, frame)} checks the keys of the frame it is given; the copy of
\code{xdpt2_p_fill} in \code{xtdpthresh.ado} is ported. \\
\code{xdpt2_is_strongly_balanced} &
message not ported &
Used to print, when \code{boottype(unit)} runs on an unbalanced effective stack, ``this is an experimental extension beyond the
common-$T$ theory'', and to store a balance flag. The package does not print that warning; the estimates are unaffected. \\
\code{xdpt2_near_common} &
message not ported &
Names, in the error message of a rank failure, the variables that are common to all units only up to rounding. The package raises
the same error with the same return code but without that hint. \\
\bottomrule
\end{longtable}

\noindent The 19 functions that were not cited before this note was written are now cited where the package implements their role
(for example \code{xdpt2_gsum_by_unit} by the unit-moment kernel, \code{xdpt2_kink_joint_V} and
\code{xdpt2_kink_joint_wind} by the joint variance, \code{xdpt2_fast_gmm_boot} by the batched bootstrap criterion,
\code{FD_con}, \code{cov_con_b} and \code{cov_con_kink_b} of \code{xthenreg} by the first differences and the Wald
covariance of the linearity test).

\input{ported_functions.tex}

\subsection*{References}

\begin{description}\small
\item[Arellano, M., and S. Bond (1991).] Some tests of specification for panel data: Monte Carlo evidence and an application to
employment equations. \emph{Review of Economic Studies} 58: 277--297.
\item[Arellano, M., and O. Bover (1995).] Another look at the instrumental variable estimation of error-components models.
\emph{Journal of Econometrics} 68: 29--51.
\item[Gong, W., and M. H. Seo (2026).] Bootstraps for dynamic panel threshold models. \emph{Journal of Econometrics} 253: 106153.
\item[Hansen, B. E. (1999a).] Threshold effects in non-dynamic panels: estimation, testing, and inference.
\emph{Journal of Econometrics} 93: 345--368.
\item[Hansen, B. E. (1999b).] The grid bootstrap and the autoregressive model. \emph{Review of Economics and Statistics} 81: 594--607.
\item[Mammen, E. (1993).] Bootstrap and wild bootstrap for high dimensional linear models. \emph{Annals of Statistics} 21: 255--285.
\item[Nguyen, D. C., and N. D. Lai (2026).] xtdpthresh: dynamic panel threshold regression for unbalanced panels, with endogenous
regressors and continuity-robust inference. SSRN Working Paper 6619058.
\item[Roodman, D. (2009).] How to do xtabond2: an introduction to difference and system GMM in Stata. \emph{Stata Journal} 9: 86--136.
\item[Seo, M. H., S. Kim, and Y.-J. Kim (2019).] Estimation of dynamic panel threshold model using Stata. \emph{Stata Journal} 19: 685--697.
\item[Seo, M. H., and Y. Shin (2016).] Dynamic panels with threshold effect and endogeneity. \emph{Journal of Econometrics} 195: 169--186.
\item[Windmeijer, F. (2005).] A finite sample correction for the variance of linear efficient two-step GMM estimators.
\emph{Journal of Econometrics} 126: 25--51.
\end{description}
